\documentclass[11pt]{article}
\usepackage[utf8]{inputenc}
\usepackage[T1]{fontenc}
\usepackage{amsmath,amssymb,amsthm}
\usepackage{graphicx}
\usepackage{booktabs}
\usepackage{listings}
\usepackage{xcolor}
\usepackage[margin=2.4cm]{geometry}
\usepackage{hyperref}

\graphicspath{{figures/}}
\DeclareMathOperator{\Tr}{Tr}
\newcommand{\C}{\mathbb{C}}
\newcommand{\M}{\mathcal{M}}
\newtheorem{proposition}{Proposition}

\lstset{language=Python, basicstyle=\ttfamily\footnotesize, frame=single, breaklines=true,
        columns=fullflexible, keepspaces=true, commentstyle=\color{gray}, showstringspaces=false}

% numbers and tables produced by scripts/paper/*.py (bash scripts/paper/run_all.sh [quick|full] [device])
\newcommand{\gen}[1]{\IfFileExists{generated/#1.tex}{\input{generated/#1.tex}}{}}
\gen{validation}\gen{search}\gen{bounds}\gen{benchmark}
\providecommand{\ValidationRows}{\multicolumn{6}{c}{\emph{run scripts/paper/validate\_known\_results.py}}\\}
\providecommand{\SearchRows}{\multicolumn{9}{c}{\emph{run scripts/paper/inequality\_search.py}}\\}
\providecommand{\BoundsRows}{\multicolumn{7}{c}{\emph{run scripts/paper/bounds\_search.py}}\\}
\providecommand{\BenchmarkRows}{\multicolumn{5}{c}{\emph{run scripts/paper/benchmark\_library.py}}\\}
\providecommand{\BenchHasGPU}{0}
\newcommand{\fig}[2]{\IfFileExists{figures/#1}{\includegraphics[width=#2]{#1}}{\fbox{\small figure #1 not generated}}}

\title{\texttt{torch\_vn\_algebra}: batched random operators, states, channels\\
and constrained optimisation in finite-dimensional Type~I von Neumann algebras}
\author{Irina Nikolaeva \and Andrej Novikov}
\date{}

\begin{document}
\maketitle

\begin{abstract}
\texttt{torch\_vn\_algebra} is a PyTorch library for numerical work in algebras
$\M=\bigoplus_{c=1}^C M_{k_c}(\C)$, the finite-dimensional Type~I von Neumann algebras that describe
quantum systems with superselection rules. Operators, states and channels are stored as batched,
block-diagonal tensors, so that thousands of random samples, or thousands of independent starts of
an optimisation, are processed at once on a CPU or a GPU. The library provides random operators with
prescribed spectra and Haar-distributed eigenvectors, the three trace functionals that are natural on
$\M$, states and Lüders conditioning, completely positive maps within and between sectors, tensor
products and partial traces, Hamiltonian and Lindblad dynamics, differentiable parametrisations with
a prescribed Michelson contrast, and safeguards that estimate run time and memory before expensive
steps. We describe the design, show how the library is used on short complete examples, validate it
against a set of exact results from random matrix theory and quantum information, and apply it to the
trace inequalities studied in our earlier work. There, constrained optimisation replaces Monte Carlo
sampling: it recovers the sharp bounds that sampling misses in all but the smallest dimensions.
\end{abstract}

\section{Introduction}

Block-diagonal operators appear whenever a symmetry or a superselection rule splits the Hilbert space
into sectors: charge or particle-number sectors, parity, total spin, or the classical part of a hybrid
classical--quantum system. Mathematically, the observables form a finite-dimensional Type~I von Neumann
algebra $\M=\bigoplus_c M_{k_c}(\C)$, whose centre is spanned by the sector projections, and the choice
of trace on $\M$ is no longer unique. The questions we had in mind when writing this library come from
operator theory: which inequalities between traces of products characterise the trace or the centre of
$\M$~\cite{gardner1979,petz1988,novikov2015,virosztek2017}, and how far do they fail in terms of a
spectral quantity such as the Michelson contrast~\cite{abed2024,novikov2020}?

General-purpose quantum software (QuTiP~\cite{johansson2013}, Qiskit~\cite{qiskit2019}, the JAX-based
\texttt{dynamiqs}~\cite{dynamiqs}) is built around a single Hilbert space; sectors, several traces and
large batches of random operators with prescribed spectra are not first-class objects there. The
niche of \texttt{torch\_vn\_algebra} is this combination: the algebra $\M$ and its traces as the basic
objects, batched random sampling and batched optimisation over them, and enough quantum information
and dynamics (states, channels, tensor products, Lindblad evolution) to formulate physical questions.
We do not aim to replace the packages above for single-system simulations.

This text supersedes the description of the first version of the library~\cite{v1}; the list of
corrections to that version is part of the repository.

\paragraph{Contributions.} (i) A batched block-diagonal data model with three trace conventions and
explicit conversions between them (Sections~\ref{sec:math},~\ref{sec:design}). (ii) States, channels
within and between sectors, tensor products and dynamics that respect the sector structure
(Section~\ref{sec:usage}). (iii) A validation suite against exact results (Section~\ref{sec:validation}).
(iv) Differentiable parametrisations with a hard Michelson-contrast constraint and batched multistart
optimisation, applied to sharp bounds of trace inequalities (Section~\ref{sec:optimisation}).
(v) Run-time and memory safeguards (Section~\ref{sec:design}).

\section{Mathematical setting}\label{sec:math}

\paragraph{Algebra and traces.} Let $\M=\bigoplus_{c=1}^C M_{k_c}(\C)$ act on
$\mathcal H=\bigoplus_c\C^{k_c}$; $A\in\M$ is block diagonal, $A=\bigoplus_c A_c$. Every trace on $\M$ is
$\operatorname{tr}_w(A)=\sum_c w_c\Tr A_c$ with weights $w_c>0$. The library implements
\[
  \Tr_{\mathrm{blunt}}:\ w_c=1,\qquad \Tr_{\mathrm{norm}}:\ w_c=\tfrac1{k_c},\qquad
  \tau_{\mathrm{vN}}:\ w_c=\tfrac1{Ck_c},
\]
where $\tau_{\mathrm{vN}}$ is the faithful normal tracial state that gives every sector the same weight.
For $A\in\M$, $|A|=(A^*A)^{1/2}$; for $A\ge 0$ the Michelson contrast is
$\Delta(A)=(\lambda_{\max}-\lambda_{\min})/(\lambda_{\max}+\lambda_{\min})\in[0,1]$, with the extreme
eigenvalues taken over all sectors; $\Delta(A)=0$ if and only if $A$ is a multiple of the identity.

\paragraph{States.} A state $\omega$ on $\M$ is represented by its density $\rho\ge0$ with respect to
$\Tr_{\mathrm{blunt}}$, $\omega(A)=\Tr(\rho A)$, $\sum_c\Tr\rho_c=1$. The restriction of $\omega$ to the
centre is the classical distribution $p_c=\Tr\rho_c$ over sectors. The density of the same state with
respect to $\operatorname{tr}_w$ is $\rho^{(w)}_c=\rho_c/w_c$; the library converts between the
conventions explicitly, because statements such as ``trace preserving'' or ``the dual map'' depend on
them as soon as the sectors change.

\paragraph{Maps between sectors.} Let $\M_{\mathrm{in}}=\bigoplus_c M_{k_c}$ and
$\M_{\mathrm{out}}=\bigoplus_d M_{m_d}$. Every normal completely positive map
$\Phi:\M_{\mathrm{in}}\to\M_{\mathrm{out}}$ (in the Schrödinger picture) is a sum of completely
positive maps between blocks,
\[
  \Phi(\rho)_d=\sum_c\sum_i K^{dc}_i\,\rho_c\,(K^{dc}_i)^*,\qquad K^{dc}_i:\C^{k_c}\to\C^{m_d}.
\]
It preserves $\Tr_{\mathrm{blunt}}$ iff $\sum_{d,i}(K_i^{dc})^*K_i^{dc}=1_c$ for all $c$, and
$T_{dc}=\Tr\Phi_{dc}(\rho_c)/\Tr\rho_c$ is the classical transition matrix between sectors. The dual
with respect to $\operatorname{tr}_w$ on both sides is
$\Phi^{*,w}(A)_c=\sum_{d,i}(w^{\mathrm{out}}_d/w^{\mathrm{in}}_c)(K_i^{dc})^*A_dK_i^{dc}$; for
sector-preserving maps on one algebra the weights cancel and all three traces give the same dual.

\paragraph{Tensor products.} $\bigl(\bigoplus_cM_{k_c}\bigr)\otimes\bigl(\bigoplus_dM_{m_d}\bigr)
=\bigoplus_{(c,d)}M_{k_cm_d}$ is again of Type~I with sectors labelled by pairs. The partial trace
$(\Tr_2X)_c=\sum_d\Tr_{m_d}X_{(c,d)}$ is a map between sectors in the sense above; its dual with
respect to $\Tr_{\mathrm{blunt}}$ is the embedding $A\mapsto A\otimes1$.

\paragraph{Dynamics.} For $H=H^*\in\M$ and jump operators $L_j\in\M$ the
Gorini--Kossakowski--Sudarshan--Lindblad equation
$\dot\rho=-i[H,\rho]+\sum_j\gamma_j(L_j\rho L_j^*-\tfrac12\{L_j^*L_j,\rho\})$ preserves every sector,
so the sector probabilities $p_c$ are conserved (superselection).

\section{Design}\label{sec:design}

\paragraph{Data layout.} An operator is a tensor of shape $(B,C,k_{\max},k_{\max})$: $B$ samples, $C$
sectors, block $c$ in the top-left $k_c\times k_c$ corner and zeros elsewhere. Linear algebra is applied
to the batch of active blocks of each sector, so one call to a batched LAPACK/cuSOLVER routine serves all
samples. The padding costs memory when the sector sizes differ strongly; a bucketed layout is planned
(Section~\ref{sec:limitations}).

\paragraph{Objects.} \texttt{Operator} (with \texttt{DensityMatrix} as a subclass for states),
\texttt{Channel} (sector preserving) and \texttt{InterSectorChannel} (between sectors or algebras) are
thin wrappers around such tensors; Table~\ref{tab:modules} lists the modules. Operations on operators
are lazy: they record how to compute the result and do so on first use; the result is cached and the
recipe dropped, so intermediates of long expressions can be freed. Laziness defers and skips work but
does not fuse operations; where this matters the library uses direct formulas, e.g.\
$\Tr(\rho A)=\sum_{ij}\rho_{ij}A_{ji}$ in $O(k^2)$ instead of forming $\rho A$.

\begin{table}[h]
\centering\small
\begin{tabular}{ll}
\toprule
module & content \\
\midrule
\texttt{algebra} & \texttt{TypeIAlgebra}, \texttt{Operator}: random operators, unitary ensembles, traces,
functional calculus\\
\texttt{states} & \texttt{DensityMatrix}, random states, Gibbs states, Lüders rule, entropies, fidelity\\
\texttt{channels} & \texttt{Channel}, \texttt{InterSectorChannel}, Choi matrices, standard and random channels\\
\texttt{composite} & tensor products, \texttt{kron}, partial traces\\
\texttt{dynamics} & exact Schrödinger/von Neumann evolution, RK4, Lindblad equation, $e^{t\mathcal L}$\\
\texttt{optimize} & parametrisations with prescribed contrast, batched multistart \texttt{extremize}\\
\texttt{cost} & run-time and memory estimates, warnings and errors\\
\bottomrule
\end{tabular}
\caption{Modules of \texttt{torch\_vn\_algebra}.}
\label{tab:modules}
\end{table}

\paragraph{Random unitaries.} Haar unitaries on $U(n)$ and $O(n)$ are obtained from the QR decomposition
of a Ginibre matrix with i.i.d.\ circular Gaussian entries after fixing the phases of
$\operatorname{diag}R$~\cite{mezzadri2007}; $SU(n)$ and $SO(n)$ by multiplying with $\det^{-1/n}$ or by
a sign flip of one column; the circular orthogonal and symplectic ensembles as $W^TW$ and $W^RW$ with
$W$ Haar on $U(n)$~\cite{mehta2004}. Section~\ref{sec:validation} checks all of them.

\paragraph{Precision and differentiation.} Algebras are created in single (default) or double
precision. All operations are composed of differentiable PyTorch primitives, so gradients with respect
to the matrix entries are available; eigenvalue and singular-value decompositions have singular
gradients at degenerate spectra, which the parametrisations of Section~\ref{sec:optimisation} avoid.

\paragraph{Safeguards.} Batched computations make it easy to request hours of work or more memory than
the device has. The \texttt{cost} module therefore (a) times the first few items of every large
batched eigen-, singular-value, QR or exponential decomposition and issues a warning with the expected
total time when it exceeds a threshold (10\,s by default), (b) compares the estimated peak memory of
large allocations with the free memory of the device, warning above 50\,\% and raising an
\texttt{InsufficientMemoryError} above 90\,\%, and (c) reports the expected duration of time-stepping
and optimisation loops after their second step, optionally with a progress bar. All thresholds can
be changed, and the checks can be switched off locally.

\section{Usage by example}\label{sec:usage}

The listings of this section are complete programs from \texttt{examples/paper}; the test suite runs
them, so they stay in sync with the code. They run on a CPU in seconds and use a GPU when available.

\paragraph{Random operators with a prescribed contrast.} Listing~\ref{lst:ex1} builds $10^4$ positive
operators in $M_2\oplus M_3\oplus M_4$ whose spectra have contrast $0.6$ and whose eigenvectors are
Haar-random in every sector, and evaluates traces and spectral functions.
\lstinputlisting[caption={Random operators (\texttt{ex1\_random\_operators.py}).},label=lst:ex1]{../../examples/paper/ex1_random_operators.py}

\paragraph{States, measurement and trace conventions.} Listing~\ref{lst:ex2} conditions random states on
the outcome of a projective measurement (Lüders rule), which moves all probability into one sector, and
checks that expectation values do not depend on the trace used to represent the state.
\lstinputlisting[caption={States and conditioning (\texttt{ex2\_states\_measurement.py}).},label=lst:ex2]{../../examples/paper/ex2_states_measurement.py}

\paragraph{A map between sectors.} Two fermionic modes have sectors $N=0,1,2$ of dimensions $1,2,1$.
Particle loss with Kraus operators $\sqrt{1-gN}$, $\sqrt g\,a_1$, $\sqrt g\,a_2$ maps sector $N$ to
$N$ and $N-1$. Listing~\ref{lst:ex3} builds it from its blocks; the classical part is the transition
matrix of a death process,
\[
  T=\begin{pmatrix}1&g&0\\0&1-g&2g\\0&0&1-2g\end{pmatrix},
\]
which the program prints. The last lines illustrate a simple fact: a unital positive map cannot
increase the Michelson contrast, since $\lambda_{\min}1\le A\le\lambda_{\max}1$ implies the same bounds
for $\Phi(A)$; the non-unital loss channel increases it for most samples.
\lstinputlisting[caption={Particle loss between sectors (\texttt{ex3\_channels.py}).},label=lst:ex3]{../../examples/paper/ex3_channels.py}

\paragraph{Open-system dynamics.} Listing~\ref{lst:ex4} integrates spontaneous emission of a two-level
atom with the RK4 solver and with the exact exponential of the Lindblad generator, and compares both with
$P_e(t)=P_e(0)e^{-\gamma t}$ and $|\rho_{eg}(t)|=|\rho_{eg}(0)|e^{-\gamma t/2}$.
\lstinputlisting[caption={Spontaneous emission (\texttt{ex4\_dynamics.py}).},label=lst:ex4]{../../examples/paper/ex4_dynamics.py}

\paragraph{Tensor products.} Listing~\ref{lst:ex5} estimates the mean entanglement entropy of random
pure states on $\C^m\otimes\C^n$ and compares it with Page's formula~\cite{page1993}.
\lstinputlisting[caption={Page entropy (\texttt{ex5\_page\_entropy.py}).},label=lst:ex5]{../../examples/paper/ex5_page_entropy.py}

\paragraph{Searching for a violation.} Listing~\ref{lst:ex6} compares Haar sampling with batched gradient
ascent for $\sup_U\bigl(|\Tr(XUY)|-\Tr(XY)\bigr)$, whose exact value $\|YX\|_1-\Tr(XY)$ is positive for
non-commuting $X,Y\ge0$. Section~\ref{sec:optimisation} studies this systematically.
\lstinputlisting[caption={Sampling versus optimisation (\texttt{ex6\_inequality\_search.py}).},label=lst:ex6]{../../examples/paper/ex6_inequality_search.py}

\section{Validation against exact results}\label{sec:validation}

Table~\ref{tab:validation} compares Monte Carlo estimates produced with the library (double precision)
with exact values: fourth moments of Haar unitaries and orthogonal matrices; the Weingarten-calculus
value $\mathbb E\,\tau\bigl((AUBU^*)^2\bigr)=-1/(N^2-1)$ for traceless involutions $A=B$ and Haar $U$;
Page's formula for $\langle S(\rho_A)\rangle$; and the mean purity $(k+r)/(kr+1)$ of induced random
states of rank $r$~\cite{zyczkowski2001}. The last column is the deviation in units of the standard
error; all are of order one. For the circular ensembles the Wigner surmise is only an approximation to
the exact spacing distribution, so no $z$-score is given; Figure~\ref{fig:spacings} shows the
distributions, with $L^1$ distances
\IfFileExists{generated/validation.tex}{\ValLoneOne, \ValLoneTwo\ and \ValLoneFour\ for
$\beta=1,2,4$}{(not generated)} from the surmise; level repulsion $P(s)\sim s^\beta$ is clearly
visible. The spontaneous-emission example agrees with the analytic solution to
\IfFileExists{generated/validation.tex}{$\ValEmissionExp$ (exact exponential) and $\ValEmissionRK$ (RK4)}{(not generated)}.

\begin{table}[h]
\centering\small
\begin{tabular}{llrrrr}
\toprule
quantity & parameters & exact & Monte Carlo & std.\ error & $z$\\
\midrule
\ValidationRows
\bottomrule
\end{tabular}
\caption{Monte Carlo estimates versus exact values\IfFileExists{generated/validation.tex}{ (\ValMode\ run on
\ValDevice, PyTorch \ValTorch)}{}.}
\label{tab:validation}
\end{table}

\begin{figure}[h]
\centering
\fig{spacings.pdf}{\textwidth}\\[2mm]
\fig{page.pdf}{0.45\textwidth}
\caption{Top: nearest-neighbour spacings of COE, CUE and CSE eigenphases generated by the library,
against the Wigner surmise. Bottom: mean entanglement entropy of random pure states against Page's
formula (error bars: two standard errors).}
\label{fig:spacings}
\end{figure}

\section{From sampling to optimisation}\label{sec:optimisation}

The first version of this work~\cite{v1} studied three trace functionals by Monte Carlo:
\begin{align*}
  \text{Exp.~1:}\quad & z=|\Tr(XUY)|-\Tr(XY), && X,Y\ge0,\ U\in\M\ \text{unitary},\\
  \text{Exp.~2:}\quad & z=\Tr|XY|-\Tr(|X|\,|Y|), && X=UX_0,\ Y=VY_0,\ X_0,Y_0\ge0,\\
  \text{Exp.~3:}\quad & z=\Tr(Y|X|Y)-\Tr|YXY|, && X=X^*,\ Y\ge0,
\end{align*}
with $\Tr=\Tr_{\mathrm{blunt}}$, operators normalised to $\lambda_{\max}=1$, and the contrasts of the
operators as parameters. Two observations motivate a different approach.

\paragraph{Sampling misses the extremes.} For fixed $X,Y$ the supremum in Exp.~1 is known exactly,
$\sup_U z=\|YX\|_1-\Tr(XY)>0$. Table~\ref{tab:search} and Figure~\ref{fig:search} show that the best of
\IfFileExists{generated/search.tex}{\SearchSamples}{many} Haar samples is already negative for $k\ge4$ and
moves further away from the supremum as $k$ grows, while batched gradient ascent
(\IfFileExists{generated/search.tex}{\SearchStarts\ starts, \SearchSteps\ steps}{multistart}) finds it
to a relative accuracy of $10^{-5}$ or better. This is concentration of measure: the unitaries close to
the maximiser form a set whose Haar measure decays rapidly with the dimension. In particular, the absence
of violations in Monte Carlo samples for $k=16$ reported in~\cite{v1} is a property of the sampling, not
of the inequality.

\begin{table}[h]
\centering\small
\begin{tabular}{rcccrccrr}
\toprule
$k$ & \multicolumn{2}{c}{violation found} & $\sup z/\Tr XY$ & best sample & \multicolumn{2}{c}{ascent rel.\ error}
& \multicolumn{2}{c}{time [s]}\\
& sampling & ascent & & $z/\Tr XY$ & median & max & sampl. & ascent\\
\midrule
\SearchRows
\bottomrule
\end{tabular}
\caption{Exp.~1 for fixed random $X,Y$: Haar sampling versus gradient ascent (medians over
\IfFileExists{generated/search.tex}{\SearchPairs}{several} pairs\IfFileExists{generated/search.tex}{;
\SearchMode\ run on \SearchDevice}{}).}
\label{tab:search}
\end{table}

\begin{figure}[h]
\centering\fig{search.pdf}{0.55\textwidth}
\caption{Best $z$ found by sampling and by gradient ascent, relative to $\Tr(XY)$, against the exact
supremum (symmetric log scale).}
\label{fig:search}
\end{figure}

\paragraph{Constrained optimisation.} The natural objects are the envelopes
\[
  \overline z(\delta)=\sup\{z:\Delta(X)=\delta\},\qquad \underline z(\delta)=\inf\{z:\Delta(X)=\delta\},
\]
(and the same with $\Delta(Y)$, or with both contrasts prescribed), the supremum being taken over all
remaining variables. The module \texttt{optimize} parametrises the constraint set exactly: the spectrum
is $\lambda_{\max}=1$, $\lambda_{\min}=(1-\delta)/(1+\delta)$ and the other eigenvalues in between
(through a sigmoid), the eigenvectors are $U=U_0e^{G-G^*}$ with a Haar-random $U_0$ per start and a chart
that is re-centred after every round, and for Exp.~3 the signs of the eigenvalues of $X$ are drawn per
start and kept fixed. Every grid point and every start is one element of the batch, so the whole envelope
is computed in a single batched optimisation.

Figures~\ref{fig:bounds} and~\ref{fig:bounds2d} and Table~\ref{tab:bounds} show the result. Every Monte
Carlo sample of~\cite{v1} lies between the computed envelopes; samples that fall outside the linear
interpolation between grid points are re-optimised at their exact contrast, after which none remains
outside. The envelopes are much wider than the sample clouds, already for $k=2$.

Two envelopes can be checked analytically. In Exp.~3, $z\ge0$ always:

\begin{proposition}
For $X=X^*$ and $Y\ge0$, $\Tr|YXY|\le\Tr(Y|X|Y)$, with equality if $Y$ commutes with $X$.
\end{proposition}
\begin{proof}
With $X=X_+-X_-$, $|X|=X_++X_-$, the triangle inequality for the trace norm gives
$\|YX_+Y-YX_-Y\|_1\le\Tr(YX_+Y)+\Tr(YX_-Y)$. If $XY=YX$, then $|YXY|=Y^2|X|$.
\end{proof}
\noindent The computed lower envelope of Exp.~3 is zero to machine precision. In Exp.~1, the lower envelope
for prescribed $\Delta(X)=\delta$ equals $-(kC-2\delta/(1+\delta))$ when all sectors have size $k\ge2$:
$\Tr(XY)\le\Tr X\le kC-1+\lambda_{\min}(X)$ with equality for $Y=1$, while $|\Tr(UYX)|$ can be made zero
by a cyclic permutation of the eigenbasis of $X_c$ in every sector. The optimiser reproduces this curve to a relative deviation of
\IfFileExists{generated/bounds.tex}{\BoundsExactDev}{(not generated)}. The upper envelopes of all three
experiments are, to our knowledge, not known in closed form; the computed curves are lower bounds for
them (a local search) and suggest conjectures for sharp contrast-dependent inequalities, which we leave
for future work.

\begin{table}[h]
\centering\small
\begin{tabular}{rrrrrcr}
\toprule
Exp. & $k$ & $C$ & $\inf z$ & $\sup z$ & Monte Carlo range of $z$ & outside\\
\midrule
\BoundsRows
\bottomrule
\end{tabular}
\caption{Global extremes of $z$ over all contrasts found by constrained optimisation, the range of the
Monte Carlo samples, and the number of samples outside the envelopes after re-optimisation
\IfFileExists{generated/bounds.tex}{(\BoundsStarts\ starts per grid point, \BoundsSteps\ steps; \BoundsMode\
run on \BoundsDevice)}{}.}
\label{tab:bounds}
\end{table}

\begin{figure}[p]
\centering
\fig{bounds_exp1_k2_C1.pdf}{0.9\textwidth}\\
\fig{bounds_exp2_k2_C1.pdf}{0.9\textwidth}\\
\fig{bounds_exp3_k2_C1.pdf}{0.9\textwidth}
\IfFileExists{figures/bounds_exp1_k16_C1.pdf}{\\\fig{bounds_exp1_k16_C1.pdf}{0.9\textwidth}}{}
\caption{Monte Carlo samples (grey) and the upper (red) and lower (blue) envelopes of $z$ as functions of
one contrast, the other one being optimised as well.}
\label{fig:bounds}
\end{figure}

\begin{figure}[h]
\centering
\fig{bounds2d_exp1_k2_C1.pdf}{0.8\textwidth}
\caption{Exp.~1, $k=2$, $C=1$: $\sup z$ and $\inf z$ with both contrasts prescribed.}
\label{fig:bounds2d}
\end{figure}

\section{Performance}\label{sec:performance}

Table~\ref{tab:benchmark} and Figure~\ref{fig:benchmark} report the wall-clock time per sample of library
operations (not of isolated kernels), in single precision, with the batch size chosen so that the batch
holds about $2^{22}$ matrix entries\IfFileExists{generated/benchmark.tex}{ (\BenchMode\ run on \BenchDevice;
CPU with \BenchCpuThreads\ threads)}{}. Expected behaviour: for single small blocks the CPU is faster, because
kernel launches dominate; the GPU gains for many sectors and for the matrix products in channels and in
the optimisation step; batched symmetric eigen- and singular-value solvers lose much of their advantage
from $k\approx64$ on, where the GPU libraries switch from batched small-matrix kernels to general ones.

\begin{table}[h]
\centering\small
\ifnum\BenchHasGPU=1
\begin{tabular}{lrrrrrr}
\toprule
operation & $k$ & $C$ & batch & GPU [$\mu$s] & CPU [$\mu$s] & speedup\\
\else
\begin{tabular}{lrrrr}
\toprule
operation & $k$ & $C$ & batch & time per sample [$\mu$s]\\
\fi
\midrule
\BenchmarkRows
\bottomrule
\end{tabular}
\caption{Time per sample of library operations.}
\label{tab:benchmark}
\end{table}

\begin{figure}[h]
\centering\fig{benchmark.pdf}{0.8\textwidth}
\caption{\ifnum\BenchHasGPU=1 Speedup CPU/GPU \else Time per sample \fi of library operations.}
\label{fig:benchmark}
\end{figure}

\section{Limitations and roadmap}\label{sec:limitations}

\begin{itemize}
\item \emph{Layout.} Blocks are padded to $k_{\max}$; for very different sector sizes (e.g.\ binomial
sectors of spin chains) a bucketed layout is needed. Merging sectors of equal total charge in tensor
products (fusion rules) is not implemented yet.
\item \emph{Optimisation.} The search is local; envelopes are lower bounds for suprema (upper bounds for
infima). Gradients of spectral functions are singular at degenerate spectra; the parametrisations avoid
this for the contrast constraint, but not for every user-defined objective.
\item \emph{Scale.} Choi matrices, superoperators and $e^{t\mathcal L}$ work with $k_c^2\times k_c^2$
matrices and are limited to blocks of a few tens; RK4 integration scales to larger blocks.
\item \emph{Scope.} Finite dimensions and Type~I only; no tensor-network or sparse methods.
\end{itemize}

\section{Availability and reproducibility}

The code (MIT licence) is available at \url{https://github.com/andre-a-no/von_neumann_type_i}.
\texttt{bash scripts/paper/run\_all.sh full cuda} regenerates every table and figure of this text on a GPU;
\texttt{bash scripts/paper/run\_all.sh} runs a reduced version on a CPU in a few minutes. The tables
state on which device and in which mode they were produced.

\begin{thebibliography}{99}
\bibitem{v1} I.~Nikolaeva, A.~Novikov, ``Finite-dimensional Type I von Neumann algebras in PyTorch: a
  GPU-accelerated framework for random block-diagonal operators'', arXiv:2606.15882v1 (2026).
\bibitem{gardner1979} L.~T.~Gardner, ``An inequality characterizes the trace'', Canad.\ J.\ Math.\ 31
  (1979), 1322--1328.
\bibitem{petz1988} D.~Petz, J.~Zem\'anek, ``Characterizations of the trace'', Linear Algebra Appl.\ 111
  (1988), 43--52.
\bibitem{novikov2015} A.~Novikov, O.~Tikhonov, ``Characterization of central elements of operator
  algebras by inequalities'', arXiv:1502.01267 (2015).
\bibitem{virosztek2017} D.~Virosztek, ``Connections between centrality and local monotonicity of certain
  functions'', J.\ Math.\ Anal.\ Appl.\ 453 (2017), 221--226.
\bibitem{abed2024} S.~A.~Abed, I.~A.~Nikolaeva, A.~A.~Novikov, ``Generalisation of Michelson contrast for
  operators and its properties'', Lobachevskii J.\ Math.\ 45 (2024), 3835--3848.
\bibitem{novikov2020} A.~Novikov, S.~A.~Abed, I.~Nikolaeva, ``Generalisation of Michelson contrast for
  operators'', arXiv:2010.10130 (2020).
\bibitem{johansson2013} J.~R.~Johansson, P.~D.~Nation, F.~Nori, ``QuTiP 2: A Python framework for the
  dynamics of open quantum systems'', Comput.\ Phys.\ Commun.\ 184 (2013), 1234--1240.
\bibitem{qiskit2019} Qiskit contributors, ``Qiskit: An open-source framework for quantum computing'', 2019,
  \url{https://qiskit.org}.
\bibitem{dynamiqs} P.~Guilmin, R.~Gautier, A.~Bocquet, E.~Genois, D.~Weiss, ``dynamiqs: an open-source
  Python library for GPU-accelerated and differentiable simulation of quantum systems'',
  \url{https://github.com/dynamiqs/dynamiqs}.
\bibitem{mezzadri2007} F.~Mezzadri, ``How to generate random matrices from the classical compact
  groups'', Notices Amer.\ Math.\ Soc.\ 54 (2007), 592--604.
\bibitem{mehta2004} M.~L.~Mehta, \emph{Random Matrices}, 3rd ed., Elsevier, 2004.
\bibitem{page1993} D.~N.~Page, ``Average entropy of a subsystem'', Phys.\ Rev.\ Lett.\ 71 (1993),
  1291--1294.
\bibitem{zyczkowski2001} K.~\.Zyczkowski, H.-J.~Sommers, ``Induced measures in the space of mixed quantum
  states'', J.\ Phys.\ A 34 (2001), 7111--7125.
\bibitem{paszke2019} A.~Paszke et al., ``PyTorch: An imperative style, high-performance deep learning
  library'', NeurIPS 2019.
\end{thebibliography}

\end{document}
